\begin{result}{G.2.3}{The May 1966 written Tate-module question and its exact scope}{Primary source checked; individual leaf date retained}
A dated written formulation of the reconstruction question occurs in Grothendieck's May 1966 letter to Tate, before the Nice address. Its express starting case concerns a $p$-divisible group whose reduction is an extension of an ind-\'etale group by a toric $p$-divisible group; the final sentence then asks about more general groups. The letter supplies a written antecedent, while its catalogue status is unpublished. It does not determine the relative date of the individual undated leaf 18:5.
\end{result}
\begin{proof}[Source and comparison maps]
The first leaf of the original typescript carries the handwritten heading ``Thil, Mai 66''. The official inventory independently identifies the letter in cote~6 as dating from May 1966 \cite[cote 6]{MontpellierInventory}. On original numbered p.~30, \S3.2, the underlined question is to
\begin{quote}
``d\'eterminer, en termes d'une telle donn\'ee, quel est le module de Tate correspondant \`a la fibre g\'en\'erique de $\Phi$''.
\end{quote}
The preceding data are a Dieudonn\'e module over Witt vectors and a lift of its filtration over a complete DVR with perfect residue field; the reference to \S3.1 retains the stated reduction qualification \cite[\S\S3.1--3.2, pp.29--30]{AuditGrothendieckTate1966}. The source's first and last leaves were inspected visually. The letter is listed as unpublished in the correspondence catalogue \cite{AuditCSGCorrespondenceCatalogue}.

Here is the actual map from this degree-one question to the cohomological comparison of Part~B. Use its $K,K_0,\kappa$ and Tate conventions, take an abelian scheme $A/\mathcal O_K$ with ordinary reduction, and set $G=A[p^\infty]$. The canonical pairing gives an isomorphism of continuous $G_K$-modules
\[
 V:=H^1_{\mathrm{\acute et}}(A_{\overline K},\Q_p)
 \xrightarrow{\sim}(V_pG)^\vee.
\]
The contravariant cohomological Dieudonn\'e object is
\[
 D=H^1_{\mathrm{cris}}(A_\kappa/W(\kappa))[1/p],
\]
with its $\sigma$-semilinear Frobenius and the filtration on $D_K$
transported from $H^1_{\dR}(A_K/K)$ by crystalline/de Rham comparison.
The comparison in B.2 identifies $D_{\mathrm{cris}}(V)$ with $D$;
its inverse gives the actual continuous Galois representation
\[
 V_pG\simeq V_{\mathrm{cris}}(D)^\vee,
 \qquad
 V_{\mathrm{cris}}(D)=
 (D\otimes_{K_0}B_{\mathrm{cris}})^{\varphi=1}
 \cap\Fil^0(D_K\otimes_KB_{\dR}).
\]
Recovering the integral lattice $T_pG$ requires the integral data of Part~B as well. Thus this abelian ordinary case is a common instance of the letter and the later smooth-proper cohomological question, with the dualization explicitly retained.

For the covariant Dieudonn\'e convention, Cartier duality gives
\[
 D_{\mathrm{cov}}(G):=D_{\mathrm{contra}}(G^\vee)
 \simeq D_{\mathrm{contra}}(G)^\vee(-1),
 \qquad
 D_{\mathrm{cris}}(V_pG)=D_{\mathrm{cov}}(G)(1).
\]
On the underlying dual module the integral operators are
\[
 F_{\mathrm{cov}}(a)=\sigma\circ a\circ V_{\mathrm{contra}},
 \qquad
 V_{\mathrm{cov}}(a)=\sigma^{-1}\circ a\circ F_{\mathrm{contra}}.
\]
They are respectively $\sigma$- and $\sigma^{-1}$-semilinear and satisfy
$F_{\mathrm{cov}}V_{\mathrm{cov}}=V_{\mathrm{cov}}F_{\mathrm{cov}}=p$.
Rationally $F_{\mathrm{cov}}=p\varphi_{\mathrm{dual}}$, since
$K_0(-1)$ has Frobenius $p\sigma$.
The filtered duality is
\[
 \Fil^rD_{\mathrm{cov},K}
 =\operatorname{ann}(\Fil^{2-r}D_{\mathrm{contra},K}).
\]
This specifies the twist and annihilator filtration instead of assigning an unstated convention to the letter's symbol $M$.

Two other earlier printed sources have also been checked. The 1968 notes of Grothendieck's November--December 1966 lectures give the crystalline/de Rham comparison for a smooth proper Witt-vector lift \cite[\S7.2, pp.352--353]{AuditGrothendieckCrystals1968}. Tate's 1967 article asks for the higher-degree Hodge--Tate decomposition \cite[p.180, Remark]{Tate1967}. The latter receives the associated-graded realization of the filtered de Rham object in B.1.4; that explicitly identifies the information retained by its formulation. The inspected passages do not themselves specify the inverse filtered-Frobenius/Galois functor displayed above.

The particular question leaf is PDF scan~5, archive-numbered p.~4. Inspection at 180 dpi revealed no date or dated reused text on that leaf. The preceding scan is a tan title cover; the following handwritten sheet is smaller, and the subsequent scan is a reused SGA typescript. Their adjacency does not establish that the typescript is the verso of the question leaf. Consequently the dossier's approximate 1969--1972 range remains dossier-level evidence \cite[cover and scans 4--7]{GrothendieckCote18Scan}. The dated 1966 question is established, without asserting that it is the earliest formulation anywhere or deciding the individual leaf's chronological position.
\end{proof}
